\section{The Geometric Origin of Entanglement and the $\sigma=5$ Capacity Limit}

In the $E_8$-Persistence framework, quantum entanglement is not a non-local correlation across spacetime, but a \textbf{projection artifact} of unprojected adjacency within the 8D substrate. While two entangled particles may be kinematically distant in the $D=4$ spacetime manifold, they remain geometrically adjacent in the orthogonal $D=4$ internal symmetry sector.

\subsubsection{Entanglement as Internal Adjacency}

Because information propagates at a maximum rate of $c_{\text{eff}} = 1$ (Paper II, Eq. 5), non-locality is structurally prohibited within the projected manifold. However, entangled states share a unified topological boundary ($\chi=2$) and a phase-locked connection via the internal interaction rank ($\sigma=5$). Measurement is the act of \textbf{Address Resolution}, forcing the shared internal state to independently project onto localized spacetime coordinates. Bell's inequalities are violated because the correlation is established in the 8D substrate prior to the 4D projection, bypassing the ``speed of light'' constraint of the 4D manifold entirely.

\subsubsection{The $\sigma=5$ Multi-partite Constraint}

The interaction rank $\sigma=5$ represents the maximum number of independent, orthogonal interaction dimensions available to a single lattice node. This imposes a strict \textbf{Topological Capacity Limit} on multi-partite entanglement in \textbf{star-topology configurations} (central hub with peripheral connections). The interaction aperture $(\sigma+1)=6$ includes the central bus node plus 5 peripheral connections.

\begin{itemize}
    \item \textbf{The Coherence Wall ($N \leq 5$):} The internal symmetry manifold has sufficient degrees of freedom to embed the state. Decoherence is driven by standard thermal noise (TOL) and follows smooth exponential decay.
    \item \textbf{Topological Aliasing ($N \geq 6$):} A 6-qubit star (1 hub + 6 leaves) requires 7 independent internal coordinates. Because $(\sigma+1)=6$, the 7th connection is forced to share an internal coordinate with an existing node. This results in \textbf{spontaneous, non-thermal decoherence} as the lattice cannot maintain seven distinct phase identities.
\end{itemize}

\textbf{Critical distinction:} This prediction applies specifically to \textbf{hub-and-spoke (star) graph states}, not linear CNOT chains. In a linear chain, no single node exceeds degree-2, and the capacity limit manifests differently (as accumulated TOL across the chain length).

\subsubsection{Correlation with Existing Benchmarks}

Current quantum hardware benchmarks exhibit unresolved anomalies consistent with the $\sigma=5$ threshold:

\begin{itemize}
    \item \textbf{Mermin Inequality Violations:} Al
sina \& Latorre (2016) \cite{alsina2016} reported a ``clear degradation'' of multipartite Bell violations specifically at the $N=5$ mark on IBM superconducting hardware, which remained poorly explained by standard gate-error models.
    \item \textbf{GHZ Fidelity Scaling:} Mooney et al. (2021) \cite{mooney2021} demonstrated GHZ states up to 27 qubits on IBM devices. While standard QFT predicts smooth exponential decay ($e^{-t/T_2}$), the observed fidelity shows \textbf{super-exponential degradation} at intermediate $N$ that persists even with quantum error mitigation.
\end{itemize}

\subsubsection{Falsification Criterion: The 6-Qubit Star Coherence Cliff}

This framework yields a rigorous test: \textbf{The 6-Qubit Star Topology Decoherence Wall.}

\textbf{Experimental Design:} Prepare GHZ-like states using a \textbf{strict star topology} (one central ``hub'' qubit with CNOTs to $N$ peripheral ``leaf'' qubits), not a linear chain. Measure fidelity vs. $N$ at fixed circuit depth and temperature.

\textbf{Standard QFT Prediction:} Fidelity decreases smoothly as $F(N) \approx F_0 e^{-\alpha N}$ due to accumulated gate errors and $T_1/T_2$ decoherence. Error mitigation (QREM, parity checking) should recover fidelity proportionally.

\textbf{E$_8$-Persistence Prediction:} 
\begin{itemize}
    \item For $N \leq 5$: Fidelity follows standard decoherence curve, mitigatable by cooling and error correction.
    \item At $N = 6$: \textbf{Discontinuous fidelity drop} ($\Delta F \sim 50\%$ or more) that \textbf{cannot be mitigated} by environmental cooling, dynamical decoupling, or gate-error correction. The hardware encounters \textbf{Topological Aliasing}—attempting to process information at a density exceeding the vacuum's interaction rank.
    \item For $N > 6$: Fidelity continues degraded, possibly with additional drops at higher $N$.
\end{itemize}

\textbf{Interpretation:} If observed, this ``coherence cliff'' constitutes direct evidence of the substrate's finite interaction rank. ``Quantum noise'' above $N=5$ in star topologies is \textbf{Topological Aliasing}—the hardware attempting to encode quantum correlations at a geometric density that exceeds the resolution of the E$_8$ vacuum itself.

\subsection{The Double-Slit and Delayed-Choice as Projection Artifacts}

The same substrate mechanism explains wave-particle duality. In the $E_8$ framework, a particle passing through a double-slit apparatus occupies a \textbf{single subspace address} with geometric support over both slit trajectories in the internal $D_4$ sector. The ``wave'' is not a physical field in spacetime, but the \textbf{phase structure of the unprojected internal address}.

\subsubsection{Which-Path Detection as Forced Address Resolution}

A detector at one slit consumes channel capacity ($\nu = 16$) to resolve the subspace address at the slit location, rather than at the screen. This eliminates the phase coherence required for interference—not by ``disturbing the particle,'' but by \textbf{preempting the projection basis}.

\subsubsection{Delayed-Choice and Quantum Erasure}

There is no retrocausality. The ``delayed choice'' determines \textbf{which basis the final projection uses}:
\begin{itemize}
    \item \textbf{Which-path preserved:} Project onto internal sector eigenstates (slit A vs. B). No interference.
    \item \textbf{Which-path erased:} Project onto superposition basis (A+B vs. A-B). Interference restored.
\end{itemize}
The ``erasure'' is a \textbf{reorientation of the projection operator}, not a physical process.

\subsubsection{The $\sigma = 5$ Limit for Multi-Path Interference}

The interaction rank $\sigma = 5$ limits coherent path multiplicity. N-slit interference should show:
\begin{itemize}
    \item Normal Fourier interference for $N \leq 5$
    \item \textbf{Catastrophic degradation at $N = 6$} due to topological aliasing
\end{itemize}
This distinguishes substrate capacity limits from standard decoherence.